% Part: turing-machines
% Chapter: machines-computations
% Section: variants

\documentclass[../../../include/open-logic-section]{subfiles}

\begin{document}

\olfileid{tur}{mac}{var}
\olsection{د ټيورينګ ماشينونو بڼې}

په حقيقت کښې د ټيورينګ ماشينونو د تعريفولو ډېرې شونې لارې شته، او
زموږ تعريف يوازې يوه لار ده۔ زموږ تعريف له ځينو اړخونو څخه تر نورو
ازاد دے۔ موږ هره متناهي الفبا منو؛ يو لا محدود تعريف ښايي يوازې د
پټې دوې نښې، $\TMstroke$ او~$\TMblank$، ومني۔ موږ ماشين ته اجازه
ورکوو چې يوه نښه پر پټه وليکي او په هماغه وخت کښې وخوځېږي؛ نور
تعريفونه يا ليکل او يا خوځېدل مني۔ موږ د سر له خوځولو پرته د نښې د
ليکلو شونتيا منو؛ نور تعريفونه د~$\TMstay$ ``لارښوونه'' نۀ لري۔ زموږ
تعريف له نورو اړخونو څخه زيات محدود دے۔ موږ فرض کړې چې پټه يوازې په
يوه لوري کښې نامتناهي ده؛ نور تعريفونه پټه کيڼ او ښي دواړو لوريو ته
نامتناهي پرېږدي۔ ان هر شمېر بېلې پټې، يا د مربعونو نامتناهي جال هم
منل کېدے شي۔ موږ د ټيورينګ ماشين د لارښوونو سټ د حالت بدلون تابع په
وسيله څرګندوو؛ نور تعريفونه د حالت بدلون اړيکه کاروي چې ماشين ته په
هر ټاکلي حالت کښې تر يوې زياتې ممکنې لارښوونې ورکوي۔

د تعريف دا وروستنۍ نرمونه ځانګړې په زړه پورې ده۔ زموږ په تعريف کښې
چې ماشين په حالت~$q$ کښې نښه~$\sigma$ لولي، $\delta(q, \sigma)$
ټاکي چې نوې نښه، حالت او د پټې د سر ځاے څه دے۔ خو که د لارښوونو سټ
د اوسنيو حالت-نښه جوړو $\tuple{q,
  \sigma}$ او نويو حالت-نښه-لوري درې‌يزونو $\tuple{q', \sigma',
  D}$ تر منځ يوه اړيکه وګڼو، نو د ټيورينګ ماشين عمل ښايي په يوازيني
ډول ټاکلے نۀ وي---د لارښوونو اړيکه هم $\tuple{q,
  \sigma, q', \sigma', D}$ او هم $\tuple{q, \sigma, q'', \sigma'', D'}$
لرلے شي۔ په دې صورت کښې موږ يو \emph{ناتعييني} ټيورينګ ماشين لرو۔
دا ماشينونه د محاسبوي پېچلتيا په تيورۍ کښې مهم رول لري۔

د ټيورينګ ماشين د درېدو دپاره هم بېلابېل دودونه شته: موږ وايو هغه
وخت درېږي چې د حالت بدلون تابع تعريف شوې نۀ وي؛ نور تعريفونه غواړي
چې ماشين په يوه ځانګړي ټاکلي درېدني حالت کښې وي۔ په~\olref[hal]{sec}
کښې مو څرګنده کړې چې د ټاکلي درېدني حالت غوښتل داسې بنديز نۀ دے چې
د ټيورينګ ماشينونو محاسبه کېدونکي څيزونه بدل کړي۔ څرنګه چې زموږ د
ټيورينګ ماشينونو پټې يوازې په يوه لوري کښې نامتناهي دي، داسې حالتونه
شته چې ټيورينګ ماشين يوه لارښوونه په سمه توګه نۀ شي عملي کولے: کله
چې تر ټولو کيڼ مربع لولي او بايد کيڼ لوري ته لاړ شي۔ زموږ د تعريف له
مخې ماشين د ``لوېدو'' پر ځاے هماغلته پاتې کېږي، خو موږ داسې تعريفولے
شو چې په دې حالت کښې ودرېږي۔ دا تعريف هم معادل دے: د هغه ټيورينګ
ماشين چلند چې له مربع~$0$ څخه کيڼ لوري ته د تلو پر هڅه درېږي، د هر
حالت بدلون $\delta(q,\TMendtape) =
\tuple{q',\sigma,\TMleft}$ په ړنګولو تقليدولے شو---بيا ماشين پر
~$\TMendtape$ د کيڼ لوري ته د تلو د هڅې پر ځاے درېږي۔\footnote{دا
\emph{بالکل} کار نۀ کوي، ځکه هېڅ څيز موږ د مربع~$0$ پرته پر نورو
مربعونو د~$\TMendtape$ له ليکلو او لوستلو څخه نۀ منع کوي
(\olref[una]{ex:mover} وګورئ)۔ د دې د مخنيوي دپاره د دې موخې لپاره
يوه دويمه نښه~$\TMendtape'$ ور زياتولے شو۔}

د عددونو د څرګندولو بېلابېلې لارې هم شته (او په دې توګه د ټيورينګ
ماشين محاسبه شوې ننوت-وت تابع هم بدلېږي): موږ يوګونې څرګندونه کاروو،
خو دوه‌ګونې څرګندونه هم کارېدے شي۔ د دې دپاره پر~$\TMblank$ او
~$\TMendtape$ سربېره دوو نښو ته اړتيا ده۔

اوس يوه په زړه پورې خبره دا ده: له دې بڼو څخه يوه هم دا نۀ بدلوي چې
کومې تابعې د ټيورينګ په معنا محاسبه کېدونکې دي۔ \emph{که يوه تابع د
يو تعريف له مخې د ټيورينګ په معنا محاسبه کېدونکې وي، نو د ټولو
تعريفونو له مخې د ټيورينګ په معنا محاسبه کېدونکې ده۔}

موږ د دې د پخلي جزئياتو ته نۀ ځو۔ يوازې يوه بېلګه وګورئ: کيڼ او ښي
دواړو لوريو ته د نامتناهي پټې، يا د څو پټو، په منلو هېڅ اضافي
محاسبوي ځواک نۀ لاس ته راوړو۔ په لنډه توګه دليل دا دے چې د يوې
يو-لوري نامتناهي پټې لرونکے ټيورينګ ماشين د څو پټو يا دوه-لوري
نامتناهي پټو تقليد کولے شي۔ مثلاً، د اضافي حالتونو او لارښوونو په
کارولو د څو پټو يا دوه-لوري نامتناهي پټې لرونکي ماشين پروګرام په
داسې پروګرام ``ژباړلے'' شو چې يوه يو-لوري نامتناهي پټه لري۔ ژباړل
شوے ماشين جفت مربعونه د پټې~$1$ مربعونو دپاره (يا د دوه-لوري
نامتناهي پټې د ``مثبتو'' مربعونو دپاره)، او طاق مربعونه د پټې~$2$
مربعونو دپاره (يا د ``منفي'' مربعونو دپاره) کارولے شي۔

\end{document}
